\section{Discussion: Scope, Limitations \& Future Work}\label{sec:discussion}

The results in the preceding sections establish a clear and, we believe, defensible claim: on this rig, real chain-coupled captures carry a learnable physical signature that a diffusion verifier can use to separate genuine recordings from substrate-altered or forged ones with near-perfect rank-order separability (frame-level AUROC $=0.9999$ against shuffled emissions and $=0.9998$ against the F-A v1 synthetic forger at training step \num{100000}). It is precisely because that headline is strong that this section is devoted to stating, in detail, what the evidence does \emph{not} show. The PolieBotics programme treats anti-overclaim as a methodological discipline rather than a disclaimer, and the limitations below are integral to the result, not caveats appended to it.

\subsection{Scope: same-rig, two-session, single-subject}\label{sec:discussion:scope}

The single most important boundary on every quantitative claim in this report is that all of it is \emph{same-rig}. The corpus comprises exactly two capture sessions recorded on one projector--camera apparatus with one human performer: D2 (\texttt{cittadel\_first\_human}, \num{5992} chain rows) and V10 (\texttt{cittadel\_v10\_mainnet\_improv}, \num{3743} chain rows), for a combined \num{9735} frames. Both sessions use the same Sony IMX540 sensor delivering \SI{8}{bit} BayerRG8, the same projector, and the same physical scene geometry. Every figure in this report --- the hero substrate maps (\Cref{fig:hero}), the per-frame decision triples (\Cref{fig:overview}), the ROC and $z$-score distributions (\Cref{fig:roc,fig:zhist,fig:ecdf}), the segmentation ablation (\Cref{fig:segmentation}), and the multi-method agreement panel (\Cref{fig:multimethod}) --- is computed on this same two-session corpus.

Within that boundary the generalization we \emph{can} demonstrate is two-fold and genuine. First, it is cross-session: a verifier trained only on D2 and evaluated on held-out V10 reaches AUROC $=1.0$ for real-correct-$E$ versus wrong/shuffled-$E$, and the symmetric V10-only $\rightarrow$ D2 ablation likewise reaches AUROC $=1.0$. Second, it is cross-protocol: D2 was recorded under the v9 chain (\texttt{ROW\_DOMAIN\_TAG} \texttt{b"TB:ROW:v9"}) and V10 under the structurally distinct v10 chain (\texttt{b"TB:ROW:v10"}, which additionally folds in a \SI{32}{byte} \texttt{ai\_payload\_root} committing AI-agent response frames), so a verifier that performs across both sessions is, by construction, performing across both protocol versions on the same rig.

What these ablations do \emph{not} establish must be stated without hedging. They do not demonstrate cross-rig generalization, nor cross-camera, cross-projector, or cross-subject generalization. The held-out blocks in the \emph{main} verifier's evaluation are drawn from the same two sessions used in training, so the headline AUROC $=1.000$ is a within-session (same-rig) measurement, separated from the training rows only by held-out blocks with a 60-frame guard band and stride decimation; the single-session ablations above are the exception, and they constitute the cross-session evidence. A reader should therefore not infer that a Phase G checkpoint trained here would verify recordings made on a different apparatus, with a different sensor or projector, or featuring a different subject. We regard cross-rig validation as the central open empirical question of the programme, not a foregone conclusion.

The small number of sessions compounds this. With $n=2$ sessions and one operator, the dataset cannot disentangle session-to-session variation from rig-to-rig or subject-to-subject variation, and there is direct evidence that part of the learned coupling signal is session-specific. One mitigating detail is worth recording: the performer's appearance differs substantially between the sessions (a different costume and mask in each --- yoga attire in D2, an improv costume in V10), so the cross-session transfer already spans a change in the subject's visual appearance. It remains one performer on one rig, and we do not count this as cross-subject evidence. While the headline real-versus-shuffled discrimination holds at AUROC $=1.0$ out-of-distribution, the harder fake-pair discrimination degrades when a single-session verifier scores the held-out session: the D2-only verifier's discrimination of F-A v1 fakes on V10 drops to AUROC $=0.8588$, and real-pair-versus-zero-$E$ coupling also falls out-of-distribution (for example, the D2-only verifier on V10 reaches AUROC $=0.7221$). These drops do not represent an attacker escape --- F-A v1 remains caught --- but they show that not all of the verifier's coupling evidence transfers cleanly across sessions, which is exactly the kind of effect that more sessions and more rigs are needed to characterize.

\subsection{The strong adaptive attacker is future work}\label{sec:discussion:adaptive}

The red-team evidence in this report establishes robustness against \emph{one} concrete learned forger: F-A v1, a 42.3M-parameter ControlNet-style editor (EditorControlNet over a ConvNeXt-Tiny source encoder) trained for \num{100000} steps to produce fake captures $C_\text{fake}$ conditioned on a target emission $E_\text{target}$. Phase G catches F-A v1 at real-versus-fake AUROC $=1.0$ at every evaluated checkpoint (5k, 25k, 70k, 100k) in both sessions, and the catch is not an artifact of the forger ignoring the target emission: by step \num{100000} F-A v1 has learned measurable target-$E$ coupling (fake-correct versus fake-shuffled AUROC rising to $0.9927$ on D2 and $0.9295$ on V10), so it is caught for failing the real-capture conditional distribution, raw-capture realism, and source leakage --- not for lacking $E$-structure. \Cref{fig:attacker_progression} shows the operational consequence: the substrate-verification excess panel stays strongly red as the attacker matures (excess red declining only modestly, $67.4\% \rightarrow 52.4\%$ across steps 5k$\rightarrow$100k), while the genuine condition holds at mean real-red $0.018\%$, indicating non-convergence of the forgery under continued training of \emph{this} attacker.

We must be emphatic about the limit of that claim. The stronger adaptive attacker, F-A v2 --- a forger trained against fresh binder surrogates with a wider 17-channel hint stack and full-resolution $C_\text{fake}$ output --- exists only as a design specification. It is \textbf{not trained and not implemented}; the trainer code does not exist. Consequently, \emph{nothing in this report establishes robustness against a strong or adaptive attacker.} The F-A v1 result should be read as evidence that the substrate signature withstood one competent, non-trivial learned forger --- not as evidence that the signature is unforgeable. The off-body finding (\Cref{fig:segmentation}) sharpens this point by localizing the discriminative signal: under proper Mask R-CNN segmentation, preserving only the off-body region (E3$_{\mathrm{seg}}$) yields AUROC $=1.000$ across all twelve session-by-condition cells, while preserving only the body (E2$_{\mathrm{seg}}$) collapses to AUROC $0.009$--$0.538$. F-A v1 smooths body content and does not perturb off-body content in a chain-coupled way; whether F-A v2 can be trained to perturb off-body regions in a chain-coupled manner that matches the target emission is an open design target, not a result. The body-region collapse itself must be stated carefully, in distribution-distance / likelihood-overspread terms \cite{nalisnick2019}: under the $\epsilon$-prediction MSE score, F-A v1 fake body content lies nearer the verifier's conditional-likelihood mode than training-distribution real samples. The body signal is not absent --- it survives at low spatial frequencies (a Gaussian-blurred residual MSE at $\sigma=4$ retains E2$_{\mathrm{seg}}$ AUROC $\approx 0.998$) but is washed out by the high-frequency-dominated $\epsilon$-MSE. We deliberately avoid the phrasing ``fakes look more real than real''; the correct statement is that the absolute-MSE body-only discrimination fails via OOD overspread, while the relative-comparison chain-coupling test (per-frame top-1 $=100\%$ across 120 frames) is orthogonal to that phenomenon and remains the load-bearing chain-coupling demonstration.

\subsection{Supervised baseline blindness to fine perturbations}\label{sec:discussion:baseline}

A second limitation concerns the granularity of the evidence the verifiers actually use. The supervised CNN baseline (Phase H) is \emph{E-aware} for coarse contrasts --- it separates target emissions from shuffled and source emissions at AUROC $=1.0$ in both sessions --- but its learned evidence is coarse: it does not resolve fine XOF-level perturbations in this training regime. Concretely, Phase H is near chance on Type 1 global single-to-few-bit flips and Type 3 spatially localized bit-flips, while full emission replacement (Type 6) is trivially separable at AUROC $=1.0$. We attribute this in part to a data-augmentation limitation rather than a fundamental granularity ceiling: a sampling audit revealed per-frame label-locking, in which the \texttt{epoch\_seed} is fixed once at dataset construction and never refreshed, permanently locking each \texttt{(session, row, idx)} tuple to a single role and yielding only $\sim$17--33 unique frames per perturbation label across 28 train-pool labels. Roughly half of all frames are never shown as a negative. Phase H is therefore best characterized as a fast supervised \emph{coarse} baseline with a known and fixable augmentation deficiency, not as a production verifier.

Crucially, the same fine-perturbation structure appears in the trained Phase G diffusion verifier itself, which is independent of the label-locking artifact. An extended sensitivity evaluation (53 conditions $\times$ 2 sessions) shows Phase G's response is graded and octave/persistence-dependent. Small-$k$ Type 1 global flips are near chance (D2 $k{=}1$ AUROC $\approx 0.503$, rising only to $\approx 0.599$ at $k=4096$), and Type 3 localized flips remain near chance even at $k=256$ ($\approx 0.505$). Sensitivity emerges with magnitude and at low spatial frequencies: low-octave Type 2 perturbations reach $\approx 0.57$--$0.59$, while gross perturbations --- low-octave (oct-0) swaps, channel swaps, and full replacement --- are perfectly separable at AUROC $=1.0$. The honest framing, which we adopt throughout, is that these systems verify \emph{optically relevant} substrate alterations: a single XOF bit-flip whose rendered emission difference $\Delta E$ falls below the optical-relevance threshold simply does not produce a capture distinguishable from the genuine one, and the right comparison for any perturbation is its response against the rendered $\Delta E$, not against the bit-count of the underlying XOF change. We therefore do not claim that Phase G detects fine, cryptographically small XOF bit-flips; that fine-grained integrity is the province of the cryptographic chain itself, whose closed feedback loop --- a substituted capture $C_t$ yields a different \texttt{bayer\_hash}, hence a different $S_{t+1}$ and a different $E_{t+1}$ --- provides tamper-evidence by construction, independent of the optical verifier.

\subsection{Emissions are recoverable from the chain hash, and why this needs a bit-exact renderer}
\label{sec:discussion:regeneration}

A defining property of the protocol is that each emission $E_t$ is a \emph{deterministic} function of the chain state $S_t$ alone: given the released chain log, any party can recompute the projected pattern without the stored tile (\Cref{sec:availability}). We had occasion to exercise this property directly. The original 2023 recording made for the project's trailer is content-addressed on IPFS, and during archival one $\sim$\SI{256}{\kibi\byte} block of one emission frame became permanently unretrievable from its pin (the host reported no providers for the block). Because the frame's filename is its chain emission seed, the emission was regenerable from the hash: feeding the \SI{32}{\byte} seed through the recorded generator reproduced the frame, and we recovered the lost block this way, cross-checking against the frame's surviving byte-exact neighbours. This is the determinism guarantee of \Cref{sec:protocol_chain} made concrete --- the recording's content survives loss of its stored bytes, because it is an extendable-output-function expansion of a \SI{32}{\byte} hash.

The episode also sharpens \emph{why} the production renderer is specified to be bit-exact and integer (\Cref{subsec:pc_render}). The 2023 trailer used an earlier, floating-point generator: the seed was BLAKE3-XOF-expanded, reshaped into a small complex array, inverse-FFT'd, and bilinearly upsampled to a \texttt{float32} tile. That pipeline is mathematically deterministic but \emph{not bit-reproducible across machines}: the FFT and resize accumulate floating-point rounding in a hardware- and library-dependent order, so the same seed yields tiles that agree only to a few units in the last place on a different GPU or numerical-library build. Empirically, regenerating a known frame on contemporary hardware --- including era-correct single-precision libraries --- reproduced only $57$--$77\%$ of values exactly, the remainder off by one to a few ULP (maximum deviation $\sim 2\times10^{-5}$ on a $[0,255]$ scale); the reconstruction is thus numerically faithful but not byte-identical, and consequently does not reproduce the original content hash. The lesson is that a verifiable recording protocol must commit to a \emph{canonical} representation. The present protocol's emission renderer is integer and bit-exact by construction (\Cref{subsec:pc_render}), so its tiles \emph{do} re-derive byte-for-byte from the chain state and can be checked against a stored render exactly (\Cref{sec:availability}); the floating-point predecessor is a cautionary counter-example that motivates that design choice rather than a limitation of the protocol as specified here.

\subsection{Privacy and consent of human-subject captures}\label{sec:discussion:privacy}

Both sessions in this corpus feature an identifiable human performer (D2 the first-human session, V10 an improvisation session), and every raw frame, rendered emission tile, NPZ RGB thumbnail, and browse preview in the release inventory therefore contains identifiable-subject content. Any public release of these artifacts is consequently a privacy and consent matter, not merely a data-hosting decision. This concern is woven into the display methodology rather than bolted on: the standing ethical-display rules enforced on every promoted artifact prohibit manual body or person suppression, require common thresholds for real and altered conditions within paired panels, forbid per-frame min/max normalization, and mandate that diagnostic (non-verifier-tier) methods be explicitly labeled. Where a body-region audit is used --- for example the body-red gate under which \num{9733} of \num{9735} frames are body-safe --- it is an audit and selection proxy only and is never used to modify display thresholds. We flag explicitly that the off-body localization result of \Cref{fig:segmentation}, which is scientifically convenient (the discriminative signal lives away from the subject), is \emph{not} a license to crop or obscure the subject in released material; the rule is consistency, not concealment.

\subsection{Future work}\label{sec:discussion:future}

Three lines of work follow directly from the limitations above, in rough priority order.

\paragraph{Strong adaptive red-team (F-A v2).} The most important next step is to actually train F-A v2 against fresh, accepted binder surrogates and evaluate it against a Phase G checkpoint that has been strictly held out --- never used for the attacker's training, validation, tuning, checkpoint selection, curriculum, or logging --- so that the final evaluation is genuinely adversarial. The off-body finding gives this red-team a concrete, pre-registrable success criterion: an F-A v2 that learns to perturb off-body regions in a chain-coupled way matching the target emission should drive the E3$_{\mathrm{seg}}$ AUROC down from $1.000$, and that drop should be pre-registered as the attack-success metric before any final evaluation. Because the F-A v1 logistic probes on diffusion features already saturate at AUROC $=1.0$ (including across-checkpoint leave-one-out), the deeper feature-level decomposition that would explain \emph{which} substrate cues the verifier relies on is deliberately deferred until F-A v2 produces a target hard enough to make that decomposition informative.

\paragraph{Cross-rig and cross-subject studies.} The scope guard of \Cref{sec:discussion:scope} is also the clearest research agenda. Collecting additional sessions on different rigs (different camera, different projector), with different subjects and different scene geometry, is the only way to convert the same-rig cross-session result into a generalization claim of practical reach. The session-specific degradation already observed in the cross-session fake-pair discrimination suggests this is non-trivial and should be measured, not assumed. Two specific calibration gaps surfaced by the data inventory should be closed before or alongside such a campaign: the measured black level is currently absent (recorded as $0$ with source ``no\_measurement\_found,'' with black-level subtraction explicitly forbidden until a measured value exists), and the empirically optimal verifier-side frame-to-chain offset has not yet been measured and has differed between training pipelines (the recording convention itself is offset $0$, $\texttt{target\_chain\_row} = \texttt{capture\_row} = t$ with $E_t = \mathrm{gen}(\mathrm{XOF}(S_t))$; \Cref{subsec:pc_offset}); both should be pinned by diagnostic before any offset-dependent or photometry-dependent cross-rig claim.

\paragraph{Toward real-time verification.} The present Phase G evaluation is deliberately heavy: per genuine-versus-altered score it runs five fixed timesteps $\times$ four noise samples ($K_\text{NOISE}=4$), i.e.\ twenty forward passes per $(C,E)$ pair, at a benchmarked \SI{329}{\milli\second} per forward pass on an A100 ($\approx \SI{6.6}{\second}$ per score at batch size 1, or $\approx \SI{1.75}{\second}$ with the four noise samples batched). The chain itself commits at roughly \SI{2.5}{\hertz} (one physical exposure per chain step), so closing the gap between verification cost and capture cadence is a concrete engineering target. Promising directions include reducing the timestep/noise budget toward the small-$t$ regime where the discriminative signal is largest (the per-timestep $\Delta_\text{wrong}$ falls as $t$ rises), distilling the diffusion verifier into a single-pass scorer, and exploiting the observation that the verifier-tier signal is frame-level and pooled (per-pixel AUROC is much lower, mean $\approx 0.74$, while pooling reaches $0.9999$), which suggests a cheaper pooled statistic may suffice for an operating-point decision. We stress that all of the AUROC values reported here reflect rank-order separability between conditions on the evaluated corpus, not a deployed fixed-threshold operating point; choosing and calibrating such an operating point for a real-time deployment is itself future work and would have to be done per rig, consistent with the scope guard above.
